\SourceLeaf{135}
auch bereits zwei Jahre früher eingetreten sein, wenn die erst 1794 eingeführte Verlängerung der Dauer der Waisenpensionen bis zu dem Alter von 20 Jahren schon damals gegolten hätte. Dieser Umstand erscheint aber in so schwerer Bedeutung, wenn man (gemäss der schon oben gemachten Bemerkung) erwägt, dass 20—30 Jahre rückwärts von jener Epoche, nemlich von 1747—1757, der Durchschnittswerth der Interessentenzahl nur 21—22 betrug, 1794—1802 hingegen 35—36.

Ueber 15 ist nun, nachher, die Zahl der Pensionen nicht eher wieder gestiegen, als Ostern 1837, und hat sich seitdem immer darüber gehalten. Die Kasse hat den vollen, alle. seit 1794 eingetretenen successiven Erhöhungen mit einschliessenden Pensionsbetrag fortgezahlt, und so wenigstens \textit{implicite} — denn ob darüber vorher Verhandlungen der Kirchendeputation Statt gefunden haben, ist mir nicht bekannt — die zweite Interpretation des zweifelhaften Punkts angenommen.

Wollen wir nun aber, mit Beiseitesetzung sprachlicher und formell juristischer Rücksichten, zwischen den beiden Interpretationen nur nach innern Gründen entscheiden, so drängt sich zuvörderst sogleich die Frage auf: Wenn es wirklich unbedenklich ist, eine erhöhte Pension auch für 16 Percipienten ungeschmälert \textit{fortbestehen} zu lassen, warum soll es denn verboten sein, sie auch bei 16 Percipienten \textit{anfangen} zu lassen? Die Gefahr, wenn eine da ist, ist ja doch in dem einen Fall gerade eben so gross wie in dem andern. Dies Argument wird noch schlagender, wenn man zu grössern Zahlen fortschreitet. Denn eine erhöhete Pension bei 17 (und wie viel mehr bei 18, 19 u. s. f.) Percipienten fortbestehen zu lassen, ist doch ganz offenbar der Kasse gefährlicher, als die Erhöhung bei 16 anfangen zu lassen, und jenes ist, wenn man die zweite Interpretation annimmt, erlaubt, dieses verboten.

Zu welchen Folgerungen eine streng consequente Durchführung der zweiten Interpretation führt, ist leicht zu übersehen.

In dem gedruckten Regulativ von 1838 wird die in Rede stehende Anordnung im §. 10 mit den Worten eingeleitet: Die Witwenpension \textit{wird} theils nach dem Bestande des Fonds der Kasse, theils nach der Zahl der Witwen \textit{bestimmt}. Zufolge der der ersten Publication des Regulativs (1833) vorgesetzten Einleitung soll dasselbe die Verpflichtungen und die Rechte der Theilnehmer feststellen, und der Beitritt eines neuen Mitgliedes geschieht durch Unterschreiben eines ihm zur Kenntnissnahme von den Pflichten und Rechten vorgelegten Exemplars. Man muss also annehmen, dass das Regulativ dieselben \textit{vollständig} enthält, und dass die Kasse, gegenüber den Mitgliedern oder deren Hinterbliebenen, keine Rechte geltend machen kann, die nicht in diesem Regulativ enthalten sind.

Nun findet sich aber in demselben auch nicht ein einziges Wort als Vorbehalt, den Pensionssatz eventuell wieder zurückgehen lassen zu dürfen, es sei denn, dass man die erste Interpretation von jener Normirung des Pensionssatzes nach der Witwenzahl, annimmt. Im entgegengesetzten Fall muss man, bei strenger Consequenz, einräumen, dass die Kasse verpflichtet sei, sämmtliche, während Stattfindens von Pensionszahlen unter 16, eingetretenen Pensionserhöhungen ungeschmälert fortzuzahlen, die Zahl der Witwen (und Waisen) möge in späterer Zeit (z. B. in Folge der überhaupt erweiterten Theilnehmerzahl) \textit{so hoch anwachsen, wie sie wolle.}

Dies ist aber geradezu ungereimt, insofern man nicht annimmt, dass das Gouvernement die Gewähr zu leisten habe, was zu beurtheilen ausser meiner Competenz liegt.

Bei jedem bestimmten Zinsfusse kann der Kapitalanwachs von 5000 Thl. die Pensionserhöhung um 10 Thl. nur für eine bestimmte Anzahl von Pensionirten decken; bei dem Zinsfuss von 3 p. C\footnote{Bei Einbringung des Vorschlages zu der in Rede stehenden Regulirung, im Julius 1794, wurde mit unzweideutigen Worten, nur auf diese Verzinsung und nicht auf 4 p. C. Rechnung gemacht. Hienach ist also die Stelle in dem Aufsatz des Herrn Universitätsraths O. vom 15. December 1844, S. 5:
\begin{quote}\raggedright
‘Bei der im J. 1794 getroffenen Bestimmung .\,.\,.\,.\,. ist wahrscheinlich eine Berechnung dahin zum\par
‘Grunde gelegt, dass die Zinsen des auf 5000 Thl. erhöhten Fonds zu 4 proc. 200 Thl. betra-\par
‘gen, dass davon \(\tfrac14\) wieder zum Capital zu schlagen, \(\tfrac34\) aber unter die Witwen zu vertheilen\par
‘seien, was dann bei 15 Witwen für jede 10 Thl. betragen würde’
\end{quote}
zu berichtigen. Von einer solchen Berechnung, von 4 proc. Zinsen und von der Zurücklegung eines Viertels derselben kommt in den bald näher zu betrachtenden Verhandlungen von 1794 \textit{gar nichts} vor.} für höchstens
